% Full proofs for Kamp 1968, Chapter IV, Theorem 1 and its three lemmas.

\subsection{Unary real-order tenses and the even/odd separation}
\label{sec:annex-p09-kamp-iv1-unary-separation}

\begin{annexitem}{annex:p9.kamp.iv1.lemma1}{thm:p09-kamp-iv1-lemma1}{The complete four-set induction}
\begin{proof}
Choose a discrete, unbounded proposition \(D\) such that either
\(A=D\) or \(A=D^c\). Then
\[
                 \mathcal F(A)=\mathcal F(D)
                     =\{\varnothing,\mathbb R,D,D^c\}.
\]
This choice handles both alternatives in \(\operatorname{DU}(A)\)
within the same proof. We first establish the necessary facts about
\(D\), and then evaluate every constructor of
\(\mathsf{TL}_1\).

\paragraph{Local gaps and nearest points.}
For each \(t\in\mathbb R\), discreteness supplies
\(z<t<y\) with no point of \(D\) in \((z,t)\) or
\((t,y)\). Therefore there are points of \(D^c\) on both
sides of every \(t\). Also, no nonempty open interval is
contained in \(D\). Indeed, given \(a<b\), choose
\(r\in(a,b)\) and a right-hand gap \((r,y)\) for \(D\).
Any point strictly between \(r\) and \(\min(y,b)\)
lies in \((a,b)\setminus D\).

Assume now \(D\ne\varnothing\). Unboundedness gives a point
of \(D\) strictly on either side of each \(t\). There is in
fact a greatest such point below \(t\). To see this, let
\[
                 d_- =\sup\{d\in D:d<t\}.
\]
The set is nonempty and bounded above, so the real supremum exists.
A left-hand gap \((z,t)\) makes \(z<t\) an upper bound
for the set, hence \(d_-\leq z<t\). Discreteness at
\(d_-\) gives \(v<d_-\) with
\((v,d_-)\cap D=\varnothing\). If \(d_-\notin D\),
every member of the set in the supremum would be at most \(v\).
That contradicts the least-upper-bound property of \(d_-\).
Thus \(d_-\in D\), and maximality gives
\((d_-,t)\cap D=\varnothing\).

There is likewise a least point of \(D\) strictly above
\(t\). Define
\[
                 d_+ =\inf\{d\in D:t<d\}.
\]
The set is nonempty and bounded below. A right-hand gap
\((t,y)\) gives \(d_+\geq y>t\). Take a gap
\((d_+,w)\) to the right of \(d_+\). If
\(d_+\notin D\), all members of the set in the infimum
would be at least \(w>d_+\), contradicting its
greatest-lower-bound property. Hence \(d_+\in D\),
and \((t,d_+)\cap D=\varnothing\). These arguments
also apply when \(t\in D\); the required witnesses remain
strictly before or strictly after the current point.

\paragraph{Both temporal constructors when \(D\ne\varnothing\).}
For \(X,Y\in\mathcal F(D)\), the complete evaluation is
\[
 S(X,Y)=U(X,Y)=
 \begin{cases}
 \varnothing,&X=\varnothing\text{ or }Y\in\{\varnothing,D\},\\
 \mathbb R,&X\ne\varnothing\text{ and }Y\in\{D^c,\mathbb R\}.
 \end{cases}
 \tag{1}
\]
The alternatives are exhaustive: \(D\) is nonempty by assumption,
and \(D^c\) is nonempty by the local-gap property, so the four
members of \(\mathcal F(D)\) are distinct.

If \(X=\varnothing\), neither operation has a witness. If
\(Y=\varnothing\), a strict witness would require a nonempty
open real interval to be empty. Density rules this out. If
\(Y=D\), it would require such an interval to be contained in
\(D\), which the local-gap argument also rules out. These
observations prove the first line for both past and future.

Suppose next \(Y=\mathbb R\) and \(X\ne\varnothing\).
For \(X=D\), the two unbounded directions supply witnesses on
both sides of every current point. For \(X=D^c\), the two
local gaps supply those witnesses. For \(X=\mathbb R\),
the points \(t-1\) and \(t+1\) suffice. Every intervening
open interval is contained in \(Y\), so both operations are
true at every \(t\).

Finally suppose \(Y=D^c\). For \(X=D\), use \(d_-\)
as the past witness and \(d_+\) as the future witness. Their
open intervals to \(t\) contain no \(D\)-point, exactly as
required. For \(X=D^c\), choose
\(s_-\in(z,t)\) and \(s_+\in(t,y)\) in the two local
gaps. Both witnesses and their intervening intervals lie in
\(D^c\). The same two witnesses work for \(X=\mathbb R\).
This proves the second line of (1) in every case, including
current points which themselves lie in \(D\).

\paragraph{The empty discrete proposition.}
If \(D=\varnothing\), then
\(\mathcal F(D)=\{\varnothing,\mathbb R\}\).
Both \(S(\mathbb R,\mathbb R)\) and
\(U(\mathbb R,\mathbb R)\) equal \(\mathbb R\), with
witnesses \(t-1\) and \(t+1\), respectively. Each of the
other three pairs has either an empty witness set or an empty
interval condition, and gives \(\varnothing\). This proves
closure without treating the four displayed choices as four distinct
sets in this case.

\paragraph{Structural induction.}
Now let every letter occurring in \(\eta\) have its value in
\(\mathcal F(D)\). An atomic formula has the required value by
this hypothesis. Negation takes complement, which exchanges
\(\varnothing\) with \(\mathbb R\) and \(D\) with
\(D^c\). Conjunction takes intersection. Intersecting with
\(\varnothing\) gives \(\varnothing\), intersecting with
\(\mathbb R\) leaves the other set unchanged, equal sets
intersect to themselves, and \(D\cap D^c=\varnothing\).
Thus the Boolean steps preserve \(\mathcal F(D)\), including
when some of its displayed choices coincide.

For \(\eta=S(\alpha,\beta)\) or
\(\eta=U(\alpha,\beta)\), the induction hypotheses put
both operand truth sets in \(\mathcal F(D)\). Equation (1),
or the separate empty case, puts the output in
\(\{\varnothing,\mathbb R\}\). This exhausts the source
grammar. Since \(\mathcal F(D)=\mathcal F(A)\), it proves
the stronger statement. Taking every atomic interpretation equal
to \(A\) gives the source's stated single-proposition result
\cite[ch.~IV, lemma~1, pp.~96--97]{Kamp1968}.
\end{proof}
\end{annexitem}

\begin{annexitem}{annex:p9.kamp.iv1.unary.corollary}{thm:p09-kamp-iv1-unary-corollary}{Passing from binary temporal syntax to an arbitrary unary tense}
\begin{proof}
Fix any \(O\in\mathfrak M_{\mathcal R}\), and choose one of
its defining formulas \(\varphi(t_0)\). The real order is a
complete linear order. Apply \cref{thm:p09-kamp-iii2} with the
one-letter allowed signature \(\{P_0\}\). It supplies a
\(\mathsf{TL}_1\)-formula \(\theta\) containing no
proposition letters other than \(q_0\), and for every
\(B\subseteq\mathbb R\) and every \(t\in\mathbb R\),
\[
 \langle\mathbb R,<,B\rangle\models\varphi[t_0\mapsto t]
       \quad\Longleftrightarrow\quad
       t\in\llbracket\theta\rrbracket_{q_0\mapsto B}.
 \tag{2}
\]
The first-order translation is the literal \(D\)-translation
proved there, with the strict witness and open interval in the
same argument order. There are no surviving auxiliary predicate
letters. Any inessential letter used to encode a truth constant is
\(q_0\), which is already allowed.

Take \(B=A\) with \(\operatorname{DU}(A)\). Every letter
occurring in \(\theta\) is then interpreted by \(A\), so
\cref{thm:p09-kamp-iv1-lemma1} puts its truth set in
\(\mathcal F(A)\). By (2) that truth set is \(O(A)\).
This proves the assertion for the initially arbitrary \(O\).
No common bound on its defining formula was used. Nor was a
uniform family over other orders assumed as an extra hypothesis:
the chosen first-order formula can be interpreted on complete
orders, and its translation is restricted back to \(\mathcal R\)
to obtain (2).

For the empty and full inputs in
\cref{def:p09-kamp-iv1-du-boundaries}, take the discrete set
\(D=\varnothing\). Its four-choice family is exactly
\(\{\varnothing,\mathbb R\}\), which proves the stated
output restriction. For the singleton comparison, \(\{0\}\)
is discrete: at zero choose any sufficiently small two-sided gap,
and at a nonzero point choose a neighborhood which avoids zero.
It is nonempty and not unbounded. Its strict-past image consists
exactly of the reals \(t>0\), since its only possible witness
is zero. That set differs from all four displayed alternatives.
Finally, every \(1/m\) is isolated among the sequence values,
while for each \(y>0\) an integer \(m\geq1\) with
\(m>1/y\) gives \(0<1/m<y\). Hence the sequence set
has no right-hand empty gap at zero. This verifies both boundary
examples using the exact source definition
\cite[ch.~IV, def.~1 and corollary, pp.~96--97]{Kamp1968}.
\end{proof}
\end{annexitem}

\begin{annexitem}{annex:p9.kamp.iv1.lemma2}{thm:p09-kamp-iv1-lemma2}{All four Boolean intersection cases and their discreteness}
\begin{proof}
Let \(N=\{1,\ldots,n\}\), interpreted as the empty set
when \(n=0\), and let
\[
             \mathcal G=\{A_I:I\subseteq N\}
                          \cup\{A_I^c:I\subseteq N\}.
\]
It contains \(\varnothing=A_{\varnothing}\),
\(\mathbb R=A_{\varnothing}^c\), and each generator.
It is closed under complement. For \(I,J\subseteq N\),
its four intersection cases are
\[
\begin{aligned}
 A_I\cap A_J&=A_{I\cap J},\\
 A_I\cap A_J^c&=A_{I\setminus J},\\
 A_I^c\cap A_J&=A_{J\setminus I},\\
 A_I^c\cap A_J^c&=A_{I\cup J}^c.
\end{aligned}
 \tag{3}
\]
For the first identity, a point lying in a generator indexed by
\(I\) and one indexed by \(J\) must lie in the same generator,
by pairwise disjointness. Its index therefore belongs to
\(I\cap J\); the converse follows by inclusion. For the
second identity, membership in \(A_I\) gives an index in
\(I\), and exclusion from \(A_J\) says that index is outside
\(J\). Conversely a point in a generator with index in
\(I\setminus J\) cannot belong to any generator indexed by
\(J\), again by disjointness. Exchanging \(I,J\) proves
the third identity. The fourth is De Morgan's law together with
\(A_I\cup A_J=A_{I\cup J}\).

Thus \(\mathcal G\) is a field containing the generators, so
minimality gives
\(\mathcal B(A_1,\ldots,A_n)\subseteq\mathcal G\).
In the other direction, every \(A_I\) is a finite union of
generators. Such unions belong to the generated field because
union is expressible by complement and intersection. Their
complements also belong. Hence \(\mathcal G\subseteq
\mathcal B(A_1,\ldots,A_n)\), proving the claimed equality.
For \(n=0\) the same computation gives exactly
\(\{\varnothing,\mathbb R\}\). Empty generators do not
affect any of the identities, and can make distinct indices
represent the same set, as the statement allows.

It remains to justify the \(\operatorname{DU}\) conclusion
which will be used in the next induction. Each finite union
\(A_I\) is discrete. For nonempty \(I\), at a fixed
\(t\) choose gaps \((z_i,t)\) and \((t,y_i)\) for each
\(A_i\), \(i\in I\). Put
\[
                     z=\max_{i\in I}z_i,
                 \qquad y=\min_{i\in I}y_i.
\]
The family is finite, so \(z<t<y\), and both resulting
open intervals avoid every \(A_i\) in the union. The empty
union has such gaps trivially. If \(A_I=\varnothing\), it
is unbounded by the source convention. If it is nonempty, some
\(A_i\), \(i\in I\), is nonempty and thus has members
strictly on both sides of every real point. Those members also
belong to \(A_I\), proving its unboundedness. Therefore each
\(A_I\) is discrete and unbounded, and each of
\(A_I,A_I^c\) satisfies \(\operatorname{DU}\). The normal
form now gives that condition for every field member
\cite[ch.~IV, lemma~2, pp.~97--98]{Kamp1968}.
\end{proof}
\end{annexitem}

\begin{annexitem}{annex:p9.kamp.iv1.lemma3}{thm:p09-kamp-iv1-lemma3}{The full induction over the unary-tense language}
\begin{proof}
Write \(\mathcal B=\mathcal B(A_1,\ldots,A_n)\), and
assume that every proposition letter occurring in the finite formula
has its valuation in \(\mathcal B\). This includes the stated
valuation of \(q_1,\ldots,q_n\) and also the permitted
extension to further field-valued letters. We induct on the
construction of \(\eta\in\mathsf{TM}_{\mathcal R}\).

For an atom, the conclusion is the valuation hypothesis. If the
conclusion holds for \(\alpha\), then
\(\llbracket\neg\alpha\rrbracket_V
 =\llbracket\alpha\rrbracket_V^c\) lies in the field.
If it holds for both \(\alpha\) and \(\beta\), then
\(\llbracket\alpha\land\beta\rrbracket_V
 =\llbracket\alpha\rrbracket_V\cap
   \llbracket\beta\rrbracket_V\) also lies there.

For \(\eta=O(\alpha)\), with any
\(O\in\mathfrak M_{\mathcal R}\), put
\(C=\llbracket\alpha\rrbracket_V\).
The induction hypothesis gives \(C\in\mathcal B\).
By \cref{thm:p09-kamp-iv1-lemma2},
\(\operatorname{DU}(C)\) holds. Therefore
\cref{thm:p09-kamp-iv1-unary-corollary} gives
\[
 \llbracket O(\alpha)\rrbracket_V
       =O(C)\in\{\varnothing,\mathbb R,C,C^c\}
                        \subseteq\mathcal B.
\]
This closes the induction. Its unary step quantifies over every
operation in the alphabet; it does not choose a particular finite
basis before the formula is considered.

For completeness, adding a truth-functional connective of any
finite arity preserves the same conclusion directly. If its input
truth sets are \(C_1,\ldots,C_m\in\mathcal B\), its output
is the union, over the finitely many Boolean input vectors making
it true, of the intersections choosing \(C_i\) for a true
coordinate and \(C_i^c\) for a false one. An empty union is
\(\varnothing\); an intersection with no coordinates is
\(\mathbb R\). All these sets belong to \(\mathcal B\).
Thus constants as well as all finite-arity sentential operations
are covered, including for the empty generator family when such
constants are admitted as primitive abbreviations.

Suppose now that \(t,t'\) have the same membership in each
\(A_i\). They then have the same membership in every
\(A_I\), and also in every \(A_I^c\). By the normal
form of \cref{thm:p09-kamp-iv1-lemma2}, this holds for every
member of \(\mathcal B\). In particular it holds for the
truth set of \(\eta\), proving the asserted indistinguishability
\cite[ch.~IV, lemma~3, p.~98]{Kamp1968}.
\end{proof}
\end{annexitem}

\begin{annexitem}{annex:p9.kamp.iv1}{thm:p09-kamp-iv1}{Both nonexpressibility contradictions and the exact truth sets}
\begin{proof}
Put \(E=\{2k:k\in\mathbb Z\}\) and
\(H=\{2k+1:k\in\mathbb Z\}\). These sets are disjoint
and each is unbounded in both directions. Each is discrete in the
source's sense at every real point. If \(t\) is an integer,
the intervals \((t-1,t)\) and \((t,t+1)\) contain no
integers. If \(t\) is not an integer, choose its consecutive
integer neighbors \(m<t<m+1\); the intervals \((m,t)\)
and \((t,m+1)\) contain no integers. These gaps also avoid
\(E\) and \(H\). Thus all hypotheses of the three lemmas
are satisfied.

\paragraph{What the unary language can see.}
The field \(\mathcal B(E,H)\) consists of the four unions
\(\varnothing,E,H,E\cup H=\mathbb Z\), and their four
complements. These are exactly the eight sets displayed in
\cref{def:p09-kamp-iv1-integer-boundaries}; they are distinct
because \(E,H,\mathbb R\setminus\mathbb Z\) are three
nonempty disjoint sets partitioning \(\mathbb R\).
They are also the field's atoms: every field member either contains
each of these whole pieces or excludes it, and each piece itself
belongs to the field.

Every noninteger point lies outside both generators. Hence
\cref{thm:p09-kamp-iv1-lemma3}, with the harmless relabelling of
its two input letters to \(q_0,q_1\), shows that every
\(\mathsf{TM}_{\mathcal R}\)-formula in those letters has
the same truth value at every noninteger point under
\(q_0\mapsto E,q_1\mapsto H\). In particular it agrees
at \(-1/2\) and \(+1/2\).

\paragraph{The two strict binary witnesses.}
At \(t=1/2\), the even point zero witnesses
\(S(E,H^c)\), since \((0,1/2)\) contains no odd
integer. At \(t=-1/2\), every even point \(s<t\)
satisfies \(s\leq-2\). The odd integer \(-1\) then
lies in \((s,-1/2)\), so the interval condition fails for
every possible witness. Thus \(S(E,H^c)\) contains
\(1/2\) and excludes \(-1/2\).

In the other direction, zero witnesses \(U(E,H^c)\) at
\(-1/2\). At \(1/2\), every later even point is at
least two, so its intervening interval contains the odd integer
one. Hence \(U(E,H^c)\) contains \(-1/2\) and excludes
\(1/2\). These are exactly the two truth-value rows in the
body statement, with the second operand interpreted by the
complement of the odd integers.

\paragraph{Why either proposed definition contradicts this.}
Suppose first that a formula \(\eta(q_0,q_1)\) in the full
unary-tense language expresses \(S\) for every pair of
real propositions. Replace every occurrence of \(q_1\) in it
by \(\neg q_1\), obtaining the finite formula
\(\eta'(q_0,q_1)\). It remains in the same language and
has no additional proposition letters. Under the valuation
\(q_0\mapsto E,q_1\mapsto H\), it satisfies
\[
 \llbracket\eta'\rrbracket_{E,H}
     =\llbracket\eta\rrbracket_{E,H^c}=S(E,H^c).
 \tag{4}
\]
Here the first equality is the actual substitution property for
this language. It follows by structural induction: a \(q_0\)
atom retains value \(E\), a replaced \(q_1\) has value
\(H^c\), negation and conjunction preserve equality of the
operand truth sets, and the unary step applies the same set
function \(O\) to equal argument sets. No bound-variable
renaming is needed in this sentential syntax. The second equality
uses the assumed definition at the input pair \((E,H^c)\).

The left side of (4), being the truth set of a unary-tense
formula under \((E,H)\), must agree at \(-1/2\) and
\(1/2\). The right side does not. This contradiction excludes
every proposed definition of \(S\).

If instead \(\eta\) expressed \(U\), exactly the same
substitution induction would give
\(\llbracket\eta'\rrbracket_{E,H}=U(E,H^c)\).
The left side again agrees at the two noninteger points, while
the right side has the opposite truth values just proved. This
is a separate contradiction excluding every definition of
\(U\). The chosen language already contains all of
\(\mathfrak M_{\mathcal R}\), so restricting that alphabet
to any subfamily cannot restore expressibility. The finite
truth-functional extension is covered by the preceding lemma.

\paragraph{The full interval descriptions and their endpoints.}
We finish by checking the stronger worked descriptions of the
two separating truth sets. For every real \(t\) there is a
unique integer \(k\) with
\(2k<t\leq2k+2\). If \(t\leq2k+1\), the even
point \(2k\) witnesses \(S(E,H^c)\): its open interval
to \(t\) contains no odd integer, including when
\(t=2k+1\). If \(t>2k+1\), any even witness strictly
below \(t\) is at most \(2k\), so its interval to
\(t\) contains the odd integer \(2k+1\). Therefore
\[
                         S(E,H^c)=
                       \bigcup_{k\in\mathbb Z}(2k,2k+1].
\]
In particular when the current point is even it lies at the
right endpoint \(2k+2\) of the preceding comparison block;
the strict witness cannot be that current point, and the odd
integer in the preceding block prevents all earlier witnesses.

For the future, choose the unique integer \(k\) with
\(2k-2\leq t<2k\). If \(t\geq2k-1\), the even
point \(2k\) is a future witness whose open interval
from \(t\) avoids \(H\), also when \(t=2k-1\).
If \(t<2k-1\), every later even witness is at least
\(2k\), so the odd integer \(2k-1\) lies in the
intervening interval and prevents it. Hence
\[
                         U(E,H^c)=
                       \bigcup_{k\in\mathbb Z}[2k-1,2k).
\]
At an even current point the same strictness excludes using
that point as the future witness. At an odd current point,
the current odd integer is outside either open interval, so
both binary formulas are true. These verifications prove every
endpoint and full-truth-set assertion in the example, as well
as the two-point separation used for the theorem
\cite[ch.~IV, thm.~1, pp.~95--99]{Kamp1968}.
\end{proof}
\end{annexitem}
